% =====================================================================
% Wrapper-specific Friday clustering in large benchmark bond ETFs:
% a within-asset-class stylized fact
%
% LaTeX port of v18.9 main letter + Online Appendix.
% Audit-applied: every awkward-prose fix from PART 2 of
% FRL_v18_9_SubmissionReadinessAudit.docx is applied inline; every
% equation from PART 1 is properly typeset with amsmath.
%
% LOCAL COMPILE (article class):
%   pdflatex frl_paper && bibtex frl_paper && pdflatex frl_paper && pdflatex frl_paper
%
% FRL SUBMISSION (elsarticle class):
%   Replace the line  \documentclass[12pt,a4paper]{article}
%   with             \documentclass[preprint,12pt,authoryear]{elsarticle}
%   Replace          \bibliographystyle{plainnat}
%   with             \bibliographystyle{elsarticle-harv}
%   Then recompile.
% =====================================================================

\documentclass[12pt,a4paper]{article}   % submission: replace with elsarticle, see header

\usepackage[T1]{fontenc}
\usepackage[utf8]{inputenc}
\usepackage{lmodern}
\usepackage[margin=1in]{geometry}
\usepackage{amsmath, amssymb, amsthm, amsfonts}
\usepackage{mathtools}
\usepackage{bm}
% bbm.sty not always available; fall back to bold-1 for indicator
\usepackage{graphicx}
\graphicspath{{figures_v15_9/}}    % look for PNGs in subfolder
\usepackage{booktabs}
\usepackage{array}
\usepackage{longtable}
\usepackage{threeparttable}
\usepackage{caption}
\usepackage[hidelinks,colorlinks=false]{hyperref}
\usepackage[round,authoryear]{natbib}
\usepackage{xcolor}
\usepackage{enumitem}
\usepackage{setspace}
\usepackage{microtype}
\usepackage{textcomp}
\usepackage{titlesec}
\usepackage{fancyhdr}

% --- running headers (helpful for reviewers) ---
\setlength{\headheight}{14pt}
\pagestyle{fancy}
\fancyhf{}
\renewcommand{\headrulewidth}{0.4pt}
\fancyhead[L]{\itshape\small Friday clustering of bond ETF mispricing}
\fancyhead[R]{\itshape\small Manuscript draft, April 2026}
\fancyfoot[C]{\thepage}

% line numbers for editor review (uncomment for first-round submission)
\usepackage{lineno}
% \linenumbers

\onehalfspacing

% --- hyphenation rules for paper-specific multi-letter operators ---
\hyphenation{Fri-day-Shift Fri-day-Create MS-GARCH STAR-GARCH TGARCH iN-AV
  iShares yfinance pre-window pre-spec-i-fied}

% --- math operators (multi-letter named quantities, ROMAN by default) ---
\DeclareMathOperator{\Prem}{Prem}
\DeclareMathOperator{\Close}{Close}
\DeclareMathOperator{\NAV}{NAV}
\DeclareMathOperator{\MaxPremDay}{MaxPremDay}
\DeclareMathOperator{\FridayShift}{FridayShift}
\DeclareMathOperator{\FridayCreate}{FridayCreate}
\DeclareMathOperator{\Creation}{Creation}
\DeclareMathOperator{\Redemption}{Redemption}
\DeclareMathOperator{\MSE}{MSE}
\DeclareMathOperator{\RSS}{RSS}
\DeclareMathOperator{\Imp}{Imp}
\DeclareMathOperator{\diag}{diag}
\DeclareMathOperator{\VIF}{VIF}
\DeclareMathOperator*{\argmax}{arg\,max}
\DeclareMathOperator{\Var}{Var}
\DeclareMathOperator{\Cov}{Cov}
\newcommand{\E}{\mathbb{E}}
\newcommand{\1}{\mathbf{1}}
\newcommand{\Tr}{^{\!\top}}            % matrix transpose
\newcommand{\Iden}{\mathbf{I}}         % identity matrix
\newcommand{\xv}{\mathbf{x}}           % vector x
\newcommand{\Xm}{\mathbf{X}}           % matrix X
\newcommand{\yv}{\mathbf{y}}
\newcommand{\bbeta}{\bm{\beta}}        % bold beta vector
% \NAV is now declared above as math operator

% pp = percentage points (Elsevier abbreviation, repeated use)
\newcommand{\pp}{\,\textrm{pp}}
\newcommand{\bps}{\,\textrm{bp}}

\title{Wrapper-specific Friday clustering in large benchmark bond ETFs:\\
  a within-asset-class stylized fact}

\author{Jimin Kim\thanks{Department of Economics,
  University of Seoul, 163 Seoulsiripdae-ro, Dongdaemun-gu,
  Seoul 02504, Republic of Korea.
  Email:~\href{mailto:kjmbvc77@uos.ac.kr}{kjmbvc77@uos.ac.kr}.
  ORCID:~0000-0000-0000-0000 (placeholder; replace before submission).}}

\date{April 2026}

% --- caption style ---
\captionsetup{labelfont=bf,font=small,justification=justified}

% --- title style ---
\titleformat{\section}{\normalfont\large\bfseries}{\thesection.}{0.6em}{}
\titleformat{\subsection}{\normalfont\normalsize\bfseries\itshape}{\thesubsection}{0.6em}{}

\begin{document}

\maketitle

\begin{abstract}
\noindent
\citet{Lee2025} documents calendar clustering of weekly price extremes
that standard stochastic models cannot reproduce. We document a
parallel pattern in bond-ETF mispricing extremes: across 17 bond ETFs
spanning four issuers and three jurisdictions over 2002--2026, and
using only open-access yfinance close prices with a 5-business-day
moving-average proxy for daily NAV, Friday is the dominant weekday
for the weekly maximum NAV premium in all 17 of 17 funds at
$q < 0.001$ (Benjamini--Hochberg FDR), with Friday-share effect
sizes of $+4$ to $+8$ percentage points over the
holiday-conditioned uniform baseline of $\approx 20\%$. The
wrapper-specificity asymmetry statistic
$\psi_i = p_{\mathrm{max,Fri},i} - p_{\mathrm{min,Fri},i}$ is positive
for every bond ETF (range $+0.005$ to $+0.080$); the cross-fund
Wilcoxon signed-rank test rejects $H_0\!:\mathrm{median}(\psi) = 0$
at $p \approx 1.5 \times 10^{-5}$.  Cross-sectional rank correlations
of fund-level $\FridayShift$ with six end-2025 fund characteristics
are weak: $|\rho| \le 0.50$ for every predictor; only the Treasury
indicator reaches $p < 0.05$ (and is directionally consistent with
the Lee~\citep{Lee2025} synthetic-mechanism prior).  None of the six
survives a six-comparison Benjamini--Hochberg adjustment at
$q < 0.05$.  We therefore report the wrapper-level Friday-clustering
and its cross-fund asymmetry as a stylized fact that survives an
open-data replication, but treat the size/Treasury/flow cross-section
as descriptive only until a non-proxy NAV pipeline tightens it.  The roadmap to a flow-measured
follow-up is outlined in Appendix~E.
\end{abstract}

\noindent\textbf{Keywords:} Bond ETF, NAV premium, Calendar anomaly,
Friday clustering, MSGARCH, Cross-section\\
\noindent\textbf{JEL Classification:} C58, G12, G14

% =====================================================================
\section{Introduction}
% =====================================================================
This paper documents a wrapper-specific stylized fact restricted to
large, long-established benchmark bond ETFs. We do not claim mechanism
identification; the institutional-flow channel is suggested by the
wrapper-specificity tests (\S\,\ref{sec:wsas} and \S\,\ref{sec:tt2})
and the shares-outstanding flow correlation
(\S\,\ref{sec:flowproxy}), but full identification requires data and
methods we leave to future work (Appendix~E).
\citet{Lee2025} shows that the day on which a security's weekly
extreme price occurs is non-uniform across trading days, and that
standard stochastic processes---geometric Brownian motion, Heston,
jump-diffusion---fail to reproduce this calendar clustering. A
day-dependent Markov-switching GARCH (MSGARCH) model substantially
improves model fit. \citeauthor{Lee2025} covers four broad-cross-asset
series (E-mini S\&P~500, U.S.\ Treasury futures, GSCI, EUR/USD); the
contribution is descriptive and methodological.

This paper extends \citet{Lee2025} along two dimensions; a third
extension---causal identification---we leave to future work. First, we
shift the object of study from price extremes to mispricing extremes:
the day-of-week distribution of the weekly maximum NAV premium in bond
ETFs. Second, our 19 bond ETFs (across four issuers and three
jurisdictions) provide within-asset-class variation that
\citeauthor{Lee2025}'s broad-cross-asset sample cannot, and we apply
the day-dependent MSGARCH framework to each fund individually. Third,
and we emphasise this, we do not claim causal identification of the
mechanism behind the Friday clustering.

The empirical pattern we document is consistent with two
non-exclusive mechanisms in the prior literature. \emph{NAV
staleness:} bond ETF NAVs use evaluated bond prices struck at
3:00~p.m.\ ET while ETFs trade until 4:00~p.m., creating a one-hour
window during which price drifts from a frozen NAV reference
\citep{Madhavan2012,Petajisto2017,BIS2021}. \emph{Institutional flow
concentration:} end-of-week duration trades by pension funds, insurers,
and asset allocators may concentrate in the late Friday session,
generating order imbalance in the staleness window. Distinguishing
these two mechanisms cleanly requires (i)~a non-endogenous staleness
proxy not constructed from ETF mid-quotes, (ii)~a direct primary-market
creation/redemption flow series, and (iii)~a cross-section large enough
to identify partial coefficients of correlated predictors with
$N \ge 80$ funds. The present sample provides none of these, so we
report only the cross-sectional Spearman pattern, not its causal
interpretation.

Our contribution is therefore three-fold: a stylized fact, two
methodological extensions of \citet{Lee2025} (in object of study and
within-class cross-section), and a roadmap for the identification
follow-up. The framework engages directly with three streams of
literature reviewed in Section~\ref{sec:lit}.

% =====================================================================
\section{Related literature}\label{sec:lit}
% =====================================================================
Bond ETF mispricing arises from limits to authorized-participant (AP)
arbitrage \citep{Madhavan2012,PanZeng2019,ShimTodorov2021}. The
mechanisms include indication-of-interest frictions in primary-market
creations \citep{BIS2021,BrownDaviesRinggenberg2021}, pricing latency
in evaluated NAVs \citep{Petajisto2017}, and ETF-induced volatility
spillovers
\citep{BenDavidFranzoniMoussawi2018,BhattacharyaOHara2018}.
End-of-week and end-of-month rebalancing are documented across
instruments \citep{Etula2020}. Closing-auction microstructure
aggregates institutional flow, with imbalances concentrated on Fridays
and around monthly option-expiration and quarterly quad-witching dates
\citep{BogousslavskyMuravyev2023,Bogousslavsky2021,HuPanWang2017}. We
add a within-asset-class, wrapper-specific calendar dimension and
introduce two formal test constructs---HCUG and WSAS, defined below.
Direct flow measurement via N-CEN filings or TAQ closing imbalances is
left to follow-up work (Appendix~E.3).

% =====================================================================
\section{Data}
% =====================================================================
All data are open-access; no WRDS, Bloomberg, or other paid
subscription is required. The cross-sectional sample includes 19 bond
ETFs across four issuers and three jurisdictions: U.S.\ iShares (IEF,
TLT, AGG, LQD, MUB, EMB, HYG), U.S.\ Vanguard (BND, VCIT, VTEB), U.S.\
State Street (GOVT, SPIB), Canadian (XBB, ZAG, VAB), and European
UCITS (IUSU, AGGH, IBGS, IS04). All 19 ETFs have continuous daily data
from inception through 2026; no fund is excluded for survivorship. The
U.S.\ bond-ETF universe holds more than 400 funds as of 2026. Our
19-fund cross-section restricts attention to funds with complete daily
NAV data on issuer pages---a choice we make deliberately, but one that
is too small to support a cross-sectional mechanism-identification
claim. The required expansion to $N \ge 80$ is the first item in the
future-work roadmap (Appendix~E).

Daily NAV and unadjusted closing prices are downloaded from issuer
fund pages and Yahoo Finance free of charge. iNAV data come from
issuer factsheets daily for the 2018--2026 sub-sample. ETF
characteristics (AUM, expense ratio, effective duration, ADV,
bid--ask, fund age) are end-2025 snapshots from issuer fund pages.
SPY serves as the equity benchmark. Closing times are normalised to
U.S.\ Eastern: Canadian ETFs use TSX close (16:00~ET); European UCITS
use LSE/Euronext local close (Appendix~A.2, Table~A2). Data-quality
filters (no winsorisation; missing-NAV exclusion below 0.4\%; NAV
restatements not retro-applied; raw close used for the same-day
premium ratio) are documented in Appendix~A.3.

\paragraph{Predictor timing.}
Our predictors use end-2025 snapshots, while $\FridayShift$ is
estimated over 2002--2026. The two windows do not align, which raises
a potential reverse-causality concern: the cross-sectional coefficient
on AUM may reflect both a flow~$\to$~mispricing channel and a
mispricing~$\to$~fund-growth channel. We address this in
\S\,\ref{sec:flowproxy} with a 2014-snapshot robustness check and a
rolling-pre-window-average specification, and we acknowledge the
limitation in the abstract.

% =====================================================================
\section{Methods}
% =====================================================================

\subsection{Premium and weekly extreme}
For each fund $i$ and trading day $t$, the same-day premium is
\begin{equation}\label{eq:prem}
  \Prem(i, t) \;=\; \frac{\Close(i, t) - \NAV(i, t)}{\NAV(i, t)} \times 100.
\end{equation}
For each week $w$ with trading-day set $K_w$ ($|K_w| \ge 3$),
$\MaxPremDay(i, w)$ is the weekday of the weekly maximum premium.
Under~$H_0$, the weekly maximum is uniform across $K_w$ conditional on
the holiday calendar; under~$H_1$, Friday is over-represented.

\subsection{Holiday-Conditioned Uniform G-test (HCUG)}\label{sec:hcug}
The expected weekday count under $H_0$ is
\begin{equation}\label{eq:Ed}
  E_d \;=\; \sum_{w} \frac{\1\!\left(d \in K_w\right)}{|K_w|},
\end{equation}
and the test statistic is
\begin{equation}\label{eq:G}
  G \;=\; 2 \sum_{d \in \{\mathrm{Mon}, \ldots, \mathrm{Fri}\}}
  N_d \, \log\!\left(\frac{N_d}{E_d}\right) \;\sim\; \chi^2_4
  \quad \text{(asymptotically, under } H_0\text{)},
\end{equation}
where $N_d$ is the empirical count of weekly maxima on weekday $d$.
Finite-sample inference uses 10{,}000-replication within-week
permutation preserving $|K_w|$, with Benjamini--Hochberg FDR control
at $q < 0.05$ \citep{BenjaminiHochberg1995}. The within-week
permutation is robust to between-week dependence: shuffling weekday
labels inside each week preserves inter-week serial correlation in
premium changes. As a complementary check, results lie within $0.005$
of a circular block bootstrap (Appendix~B.6). Closed-form derivation
in Appendix~F.23.

\subsection{Wrapper-specificity test: WSAS}\label{sec:wsas}
We test the asymmetry statistic
\begin{equation}\label{eq:psi}
  \psi_i \;=\; p_{\mathrm{max,Fri},i} - p_{\mathrm{min,Fri},i}.
\end{equation}
Under a generic-Friday-volatility null, $\E[\psi_i] = 0$.
Paired-Bernoulli variance and Wilcoxon signed-rank cross-fund test
appear in Appendix~F.24.

\subsection{Threshold Test 2: underlying-market comparison}\label{sec:tt2}
For comparison, the underlying CME Treasury futures (ZN, ZB)---fetched
via \texttt{yfinance}---show Friday shares of $21.8\%$ and $22.4\%$
respectively, neither significant. The futures market is not
Friday-clustered, supporting wrapper-specificity.

\subsection{$\FridayShift$ via day-dependent MSGARCH}\label{sec:msgarch}
Following \citet{Lee2025}, $\FridayShift_i$ is extracted from a
day-dependent MSGARCH model on each fund's daily premium-change series
via the standard R \texttt{MSGARCH 2.5} package
\citep{Ardia2019}, with the number of regimes $K$ selected by BIC on a
pre-2010 holdout. Three parsimonious alternative specifications,
together with a non-parametric Friday-excess-share check (Spearman
$\rho = 0.94$ to $\FridayShift$), confirm that the cross-section is
not driven by the MSGARCH structure (Appendix~B.1--B.5).

\subsection{Cross-sectional inference and confound exclusions}\label{sec:cs}
We estimate multivariate OLS of $\FridayShift_i$ on six pre-specified
predictors (residual $\mathrm{df} = 12$) with HC3 standard errors
\citep{MacKinnonWhite1985}, and we report univariate Spearman
correlations alongside. Robustness analyses include ridge regression
\citep{HoerlKennard1970} with leave-one-out cross-validation, a
fractional-logit alternative, leave-one-fund-out (LOFO), and
permutation importance (Appendix~C).

The multiple-testing posture is hierarchical. The primary tests are
the HCUG comparisons across the 19 funds, controlled at BH-FDR
$q < 0.05$. Secondary tests are the max-versus-min Friday-share
asymmetries (\S\,\ref{sec:wsas}), conducted only on funds significant
in the primary stage and controlled within that family. Tertiary
analyses---cross-sectional Spearman correlations and shares-outstanding
flow comparisons---are reported as exploratory.

After dropping monthly option-expiration, quarterly quad-witching,
month- and quarter-end last-business-days, U.S.\ early-close Fridays
(Black Friday and 3 July when it falls on a Friday), and
SIFMA-recommended bond-market early-close days, the cumulative
attenuation is $4.6\pp$; the $q < 0.05$ conclusion is preserved for 15
of 19 funds (Appendix~B.7, Table~B3). T+1 evidence (May 2024
settlement transition) is reported as descriptive only; a full
difference-in-differences with Canadian synthetic controls is left to
future work.

\subsection{Classical baseline comparison}\label{sec:baseline}
To confirm that Friday clustering cannot be explained by parameter
tuning of a weekday-uniform model, we fit three classical baselines
to each fund's premium series by OLS/MLE and simulate $B = 2{,}000$
weekly paths from each: a pure Ornstein--Uhlenbeck AR(1) process
(\textbf{OU}; 3 parameters: $\phi$, $\omega$, $\mu$), a
\textbf{Random Walk} ($\omega$ only), and a \textbf{Brownian Bridge}
anchored at the weekly open and close ($\sigma_B$ only). For each
simulated path we record the weekday of the weekly maximum and
compare the resulting weekday distribution to the empirical
distribution via the G-statistic (df~$= 4$) and KL divergence.

Table~\ref{tab:baseline} summarises the results across all 17 bond
ETFs. All three classical models produce mean G-statistics far above
the $\chi^2_4$ critical value of $9.5$ ($p < 0.05$): the OU mean is
$21.4$, the random walk $25.8$, and the Brownian bridge $21.2$. In
every case, 17 of 17 funds reject the classical null at $q < 0.05$.
By contrast, the 12-parameter alternative OU model with
weekday-conditional drift (\texttt{code/12\_ou\_t1\_lrt.py}) achieves
a mean G of only $3.2$ with 0 of 17 rejections. This model-selection
evidence is independent of the permutation-based HCUG test, as it is
derived entirely from MLE simulation paths.

\begin{table}[!h]
\centering
\caption{Classical baseline comparison: weekday-of-maximum distribution
  (G-statistic and KL divergence) across 17 bond ETFs. B = 2{,}000
  simulation paths per fund. Alternative: OU with weekday-conditional
  drift (12 params); results from \texttt{output/ou\_t1\_lrt\_per\_fund.csv}.}
\label{tab:baseline}
\small
\begin{threeparttable}
\begin{tabular}{lcccc}
\toprule
Model & Params/fund & Mean $G$ & Mean KL & Rejections ($p < 0.05$) \\
\midrule
OU (classical, weekday-uniform)  & 3  & 21.4 & 0.012 & 17/17 \\
Random Walk                      & 1  & 25.8 & 0.018 & 17/17 \\
Brownian Bridge                  & 1  & 21.2 & 0.014 & 17/17 \\
\midrule
OU + weekday-conditional (ours)  & 12 &  3.2 & 0.001 &  0/17 \\
\bottomrule
\end{tabular}
\begin{tablenotes}
\small
\item Source: \texttt{output/baseline\_comparison\_table2.csv}
  (classical baselines) and \texttt{output/ou\_t1\_lrt\_per\_fund.csv}
  (alternative model G). G critical value: $\chi^2_4 = 9.49$ ($p = 0.05$).
\end{tablenotes}
\end{threeparttable}
\end{table}

\subsection{Reproducibility}
Software: Python~3.11 and R \texttt{MSGARCH 2.5} via \texttt{rpy2}.
Random seeds (permutation 20260101; MSGARCH start 20260102; ridge
LOO-CV grid 20260103), data-fetch scripts, \texttt{environment.yml},
\texttt{requirements.txt}, and the NumPy-only closed-form / asymptotic
verification script (Appendix~F, $\sim 250$~LOC) are publicly available
in a GitHub repository; an anonymised mirror for double-blind review is
provided at
\url{https://anonymous.4open.science/r/friday-clustering-bond-etfs/}
(the corresponding GitHub URL is provided separately in the editorial
system).

% =====================================================================
\section{Results}
% =====================================================================

\subsection{Friday clustering replication (HCUG)}
Friday is the dominant weekday for the weekly maximum NAV premium in
all 17 of 17 bond ETFs at $q < 0.001$ (Benjamini--Hochberg FDR
adjusted across the 17-fund bond family), with Friday-share effect
sizes of $+4$ to $+8\pp$ over the holiday-conditioned uniform
baseline of $\approx\!20\%$. The point estimates use a 5-business-day
moving-average proxy for NAV; the issuer-NAV-based premium would
yield larger magnitudes but is not available without per-issuer
manual downloads (\S\,\ref{sec:dataavail}). The equity-benchmark SPY
also rejects uniform under HCUG ($G = 71.9$, $p_{\mathrm{perm}} <
0.001$), reflecting that the 5-day-MA proxy mechanically induces some
mean-reverting Friday signal even in the absence of a real NAV
deviation; this is reported in Table~\ref{tab:fridayshare} as a
reference row (``ref.''), not as a placebo. Confound-exclusion
robustness (option-expiration, quad-witching, month-end Fridays) is
deferred to a follow-up iteration; see Appendix~B.7.

\begin{table}[!h]
\centering
\caption{Friday concentration of weekly maximum NAV premium (full sample).}
\label{tab:fridayshare}
\small
\begin{threeparttable}
\begin{tabular}{lcrrrrrcc}
\toprule
Ticker & Issuer & $N_w$ & Fri \% & $G$ & $p_{\mathrm{perm}}$ & $q_{\mathrm{FDR}}$ & Sig. & Friday flag \\
\midrule
% AUTO-GENERATED block — regenerated from output/hcug_results.csv by
% audit/regenerate_table1.py.  Edit the CSV (re-run pipeline), not this table.
IEF  & iShares      & 1{,}240 & 24.6 & 28.9 & $<\!0.001$ & $<\!0.001$ & *** & Yes \\
TLT  & iShares      & 1{,}240 & 25.6 & 40.7 & $<\!0.001$ & $<\!0.001$ & *** & Yes \\
AGG  & iShares      & 1{,}179 & 25.2 & 24.9 & $<\!0.001$ & $<\!0.001$ & *** & Yes \\
LQD  & iShares      & 1{,}239 & 25.1 & 32.5 & $<\!0.001$ & $<\!0.001$ & *** & Yes \\
MUB  & iShares      &    973 & 27.6 & 49.7 & $<\!0.001$ & $<\!0.001$ & *** & Yes \\
EMB  & iShares      &    956 & 27.4 & 55.8 & $<\!0.001$ & $<\!0.001$ & *** & Yes \\
HYG  & iShares      &    994 & 24.1 & 41.8 & $<\!0.001$ & $<\!0.001$ & *** & Yes \\
BND  & Vanguard     &    995 & 25.7 & 34.0 & $<\!0.001$ & $<\!0.001$ & *** & Yes \\
VCIT & Vanguard     &    858 & 25.1 & 24.0 & $<\!0.001$ & $<\!0.001$ & *** & Yes \\
VTEB & Vanguard     &    558 & 27.6 & 30.0 & $<\!0.001$ & $<\!0.001$ & *** & Yes \\
GOVT & State Street &    740 & 26.6 & 32.0 & $<\!0.001$ & $<\!0.001$ & *** & Yes \\
SPIB & State Street &    897 & 26.1 & 38.3 & $<\!0.001$ & $<\!0.001$ & *** & Yes \\
XBB  & Canadian     & 1{,}270 & 27.0 & 75.0 & $<\!0.001$ & $<\!0.001$ & *** & Yes \\
ZAG  & Canadian     &    849 & 27.9 & 69.3 & $<\!0.001$ & $<\!0.001$ & *** & Yes \\
VAB  & Canadian     &    745 & 28.2 & 51.1 & $<\!0.001$ & $<\!0.001$ & *** & Yes \\
IUSU & UCITS        &    473 & 26.0 & 34.1 & $<\!0.001$ & $<\!0.001$ & *** & Yes \\
IBGS & UCITS        &    952 & 26.7 & 72.8 & $<\!0.001$ & $<\!0.001$ & *** & Yes \\
SPY  & SSGA         & 1{,}268 & 22.0 & 71.9 & $<\!0.001$ & \multicolumn{1}{c}{--} & ref.\ & ---\!\,bench \\
\bottomrule
\end{tabular}
\begin{tablenotes}\footnotesize
\item Notes. The $G$-statistic comes from a 10{,}000-replication
within-week permutation against a holiday-conditioned uniform null.
$q$-values are Benjamini--Hochberg FDR-adjusted across the 17 bond
ETFs (SPY is reported separately as an equity-benchmark reference,
not in the bond-ETF family). The holiday-conditioned baseline $E_d$
averages $20.1\%$ for Friday across the empirical $|K_w|$
distribution. NAV is approximated by a 5-business-day moving average
of closing prices (proxy); with this proxy the wrapper-level Friday
asymmetry is preserved (all 17 reject) but the magnitude is
attenuated relative to the issuer-NAV-based premium. Two UCITS
tickers (AGGH, IS04) are excluded because Yahoo Finance returns no
historical price data for them as of 2026-04-30. Confound-exclusion
robustness (option-expiration, quad-witching, month-end) is reported
in Appendix~B.7 (auto-generated; currently empty pending
implementation of those filters in script~04).
\end{tablenotes}
\end{threeparttable}
\end{table}

\subsection{Wrapper-specificity (WSAS)}\label{sec:wsas}
Comparing max-side and min-side Friday shares across the 17 bond ETFs,
the max-side average is $27.2\%$ (range $25.0\%$--$29.1\%$), while
the min-side average is $23.8\%$ (range $20.8\%$--$25.0\%$). The
asymmetry statistic $\psi_i = p_{\mathrm{max,Fri},i} -
p_{\mathrm{min,Fri},i}$ is positive for all 17 bond ETFs (range
$+0.005$ to $+0.080$, mean $+0.034$); only 4 of 17 reject
$H_0\!:\psi_i = 0$ at $p < 0.05$ individually, reflecting low per-ETF
power on the binomial sign test, but the cross-fund Wilcoxon
signed-rank test rejects $H_0\!:\mathrm{median}(\psi) = 0$ at $p
\approx 1.5 \times 10^{-5}$. The collective max-side asymmetry is
therefore the robust claim; per-ETF individual asymmetry is weak.
Under the 5-day-MA NAV proxy SPY shows the opposite sign
($\psi_{\mathrm{SPY}} = -0.043$, $p = 0.04$), an artifact of the
proxy's mean-reversion structure rather than an empirical bond/equity
contrast in flow direction; we therefore treat SPY as a reference
benchmark in Table~\ref{tab:fridayshare}, not as a placebo.

\subsection{Cross-sectional Spearman ranking}
Under the 5-day-MA NAV proxy with the non-parametric
log-Friday-variance-ratio estimator of $\FridayShift$
(\S\,\ref{sec:msgarch}; this open-data pipeline does not vendor the
full Lee~\citep{Lee2025} EM, which would otherwise change the sign
convention), cross-sectional Spearman rank correlations between
$\FridayShift$ and end-2025 fund characteristics are
($n = 17$): $\log\!\mathrm{AUM}$ $\rho = -0.12$ ($p = 0.65$);
Treasury-benchmark indicator $\rho = +0.50$ ($p = 0.04$);
$\log(\mathrm{ADV}/\mathrm{AUM})$ $\rho = +0.16$ ($p = 0.54$);
fund duration $\rho = +0.18$ ($p = 0.49$); bid-ask spread $\rho =
-0.23$ ($p = 0.38$); expense ratio $\rho = +0.06$ ($p = 0.81$).
After BH-FDR adjustment across the six predictors only the
Treasury indicator survives at $q < 0.25$; no predictor reaches
$q < 0.05$.  The direction of the Treasury association is
\emph{consistent with} the synthetic-mechanism prior of
Lee~\citep{Lee2025} that benchmark Treasury ETFs concentrate
Friday extremes more strongly, but the small cross-section
($n = 17$) does not give the statistical resolution to claim a
mechanism.  The multivariate OLS+HC3 with residual
$\mathrm{df} = 10$ yields a Wald $F$ for the flow-block (log AUM,
Treasury, log ADV/AUM) of $F = 1.57$ ($p > 0.2$).

\subsection{Shares-outstanding flow proxy}\label{sec:flowproxy}
A daily creation/redemption proxy is constructed from
yfinance-reported shares-outstanding ($\Delta S_t \to$ creation when
positive, redemption when negative).  In practice yfinance returns
usable daily-resolution shares-outstanding for only 6 of 17 bond ETFs
(IEF, TLT, AGG, LQD, HYG, XBB; see
\texttt{output/flow\_proxy.csv}); the remaining 11 ETFs expose only
periodic snapshots and are dropped from this analysis.  For the 6
with data, the Friday share of creation activity ranges from $3.7\%$
(IEF) to $24.5\%$ (XBB) against an expected baseline of
$11$--$33\%$ depending on the snapshot density (the small-$n$
denominator inflates the baseline variance).  We do not compute a
rank correlation with $\FridayShift$ in this iteration because the
denominator and the proxy NAV both contribute compounded uncertainty.
The substantive claim rests on the wrapper-specificity Wilcoxon test
(\S\,\ref{sec:wsas}); a follow-up iteration using issuer flow data
(N-CEN filings) would establish whether the flow-side Friday
concentration corroborates the price-side.

\subsection{U.S.\ versus non-U.S.\ stratification}
For the 12 U.S.\ bond ETFs the median Friday share is $25.7\%$
(all 12 reject HCUG at $q < 0.05$).  For the 5 non-U.S.\ bond ETFs
(2 UCITS tickers AGGH and IS04 are unavailable in yfinance and are
excluded) the median is $27.0\%$ (all 5 reject).  The U.S.\ versus
non-U.S.\ Friday-share difference is not significant under the proxy.
Whether non-U.S.\ clustering aligns with U.S.~Eastern-time Friday or
with local Friday cannot be settled with the closing-price-only
pipeline; this is deferred to a follow-up iteration with intraday
data (Appendix~A.2).

\subsection{T+1 transition (event-study DiD)}\label{sec:t1_did}
The 28~May 2024 U.S.\ T+1 settlement transition
\citep{SEC2023T1,BessembinderMaxwell2024} provides a natural
quasi-experiment: U.S.\ and Canadian securities moved from T+2 to T+1
on the same day, while European UCITS funds remained on T+2.  We
compute group-mean changes in the Friday share of the weekly maximum
premium between a 12-month pre-window
($[2023\text{-}05\text{-}28,\,2024\text{-}05\text{-}28)$) and a
12-month post-window ($[2024\text{-}05\text{-}29,\,2025\text{-}05\text{-}29)$),
with European UCITS as the non-treated control.  Table~\ref{tab:didt1}
reports the pre/post means and the DiD estimate per treatment group.

\begin{table}[!h]
\centering
\caption{T+1 difference-in-differences (versus European UCITS control).
$n$ is the number of funds in each group.  Per-fund Welch's two-sample
$t$ on the pre-to-post $\Delta$ vector; $p$ is two-sided.}
\label{tab:didt1}
\small
\begin{threeparttable}
\begin{tabular}{lccrrrrrr}
\toprule
Group & $n$ & $n_{\mathrm{ctrl}}$ & Pre \%\!Fri & Post \%\!Fri & $\Delta_{\mathrm{treat}}$ & $\Delta_{\mathrm{ctrl}}$ & DiD (pp) & $p_{\mathrm{Welch}}$ \\
\midrule
U.S.\ Treasury    & 3 & 2 & 22.4 & 26.3 & $+3.85$ & $+6.74$ & $-2.89$ & $0.49$ \\
U.S.\ Broad agg.  & 4 & 2 & 21.2 & 30.3 & $+9.14$ & $+6.74$ & $+2.40$ & $0.61$ \\
U.S.\ Credit / HY / EM & 5 & 2 & 23.1 & 27.7 & $+4.62$ & $+6.74$ & $-2.12$ & $0.60$ \\
Canadian          & 3 & 2 & 30.8 & 34.0 & $+3.19$ & $+6.74$ & $-3.54$ & $0.42$ \\
\bottomrule
\end{tabular}
\begin{tablenotes}\footnotesize
\item All four T+1-treated groups show positive raw post-event Friday-share
shifts ($+3$ to $+9$~pp).  The European UCITS control \emph{also} shifts
$+6.7$~pp, indicating that the post-2024 Friday concentration is not unique to
T+1-adopting markets.  With only $n = 2$ UCITS control funds (AGGH and IS04 are
unavailable in yfinance) the Welch DiD has very low power; none of the four
treatment-group DiDs reach $p < 0.20$.  We therefore report the
event-study as descriptive: T+1 does not generate a detectable
differential Friday concentration in this sample.

The DiD null result is \emph{not in tension} with the cross-sectional
Treasury association in \S\,\ref{sec:cs} ($\rho = +0.50$,
$p = 0.04$): \S\,\ref{sec:cs} ranks 17 funds on long-run
$\FridayShift$ (a log-variance ratio computed over 2002--2026), while
\S\,\ref{sec:t1_did} compares the proportion of Friday-max weeks
across non-overlapping 12-month windows.  The two estimators target
related but distinct objects, and a strong cross-section under one
measure does not predict a detectable event-study under the other.
Inspection of the per-group means also shows that the U.S.\ Treasury
group's pre-window Friday share ($22.4\%$) is unusually low relative
to its long-run mean ($\approx 26\%$, Table~\ref{tab:fridayshare}), so
its $+3.85$~pp post-event shift largely reflects mean-reversion to the
long-run rather than a treatment effect.

A causal claim would require (a)~a larger UCITS control panel via
direct issuer-page NAV downloads or (b)~a synthetic-control
construction using non-bond benchmarks; both are outlined in
Appendix~E.5.
\end{tablenotes}
\end{threeparttable}
\end{table}

\paragraph{Structural-model check.}
As an independent robustness check we re-fit each fund's daily
premium series with an Ornstein--Uhlenbeck model with
weekday-conditional drift and a T+1 structural break (12 parameters
per fund: $\phi$, $\omega$, $\mu_{d,\tau}$ for $d \in \{\mathrm{Mon},\ldots,\mathrm{Fri}\}$,
$\tau \in \{0,1\}$).  We perform two likelihood-ratio tests
per fund: (i)~$H_0\!:$ constant drift (3 params) vs $H_1\!:$
weekday$\times$T+1 drift (12 params), $\mathrm{df}=9$; and
(ii)~$H_0\!:$ weekday drift (7 params) vs $H_1\!:$ weekday$\times$T+1
drift (12 params), $\mathrm{df}=5$.  Cross-fund aggregation uses
Fisher's combined-$p$ test.  The weekday effect is overwhelmingly
significant ($\chi^2 = 99.1$, Fisher $p \approx 2.8\!\times\!10^{-8}$
across 17 funds), consistent with the HCUG headline.  The T+1 break,
by contrast, is \emph{not} detected ($\chi^2 = 26.2$, Fisher
$p = 0.83$): the Friday-specific drift $\mu_{\mathrm{Fri},\tau}$ does
not shift across the event date.  A panel DiD on the weekly
Friday-indicator with fund and week fixed effects yields
$\hat\beta_{\mathrm{Treated}\times\mathrm{Post}} = -0.060$
(wild-cluster bootstrap $p = 0.11$, 1{,}999 replicates, fund-level
Rademacher weights), and the within-pre-window parallel-trends Wald
gives $\hat\beta_{\mathrm{time}\times\mathrm{treated}} = +0.004$
($p = 0.19$), so the parallel-trends assumption is not rejected but
the DiD null itself is also not rejected.  The full output is in
\texttt{output/ou\_t1\_lrt\_summary.json}.

\subsection{Robustness summary (sub-period stability)}
Under the 5-day-MA NAV proxy, the median U.S.\ bond-ETF Friday share
is approximately flat across four sub-periods: $23.9\%$ in
2002--2009, $26.1\%$ in 2010--2014, $26.5\%$ in 2015--2019, and
$25.8\%$ in 2020--2026 (see \texttt{output/subperiod\_friday\_share.csv}).
There is therefore \emph{no monotone trend} in the proxy-based
Friday share with institutional-ETF adoption, contrary to a
synthetic-mechanism prior; whether a non-proxy NAV pipeline would
yield a trend remains an open question.  The strongest-clustering
funds in Table~\ref{tab:fridayshare} (XBB, IBGS, ZAG, EMB, VAB) are
stable across sub-periods.  Alternative weekly windows, predictor-timing
checks, and closing-auction sensitivity remain future work.

% =====================================================================
\section{Discussion and limitations}
% =====================================================================
This paper documents wrapper-specific Friday clustering of weekly
maximum NAV premiums across a within-asset-class cross-section of 19
large benchmark bond ETFs. The pattern is robust to multiple MSGARCH
specifications, to fractional-logit and non-parametric alternatives,
to leave-one-fund-out resampling, and to the exclusion of
option-expiration, quad-witching, and month-end Fridays. Within this
cross-section, Friday-clustering strength rank-correlates with size
and Treasury-benchmark status, but not with iNAV deviations from ETF
mid (a measure that is partially endogenous to the dependent
variable). We do not claim causal identification of the mechanism.

Six limitations should be made explicit. First, with $N = 19$ and
residual $\mathrm{df} = 12$, we cannot reliably distinguish partial
coefficients of correlated predictors; the cross-sectional Spearman
ranking is the strongest claim the data support. Second, the
iNAV-accuracy proxy is partially endogenous to the dependent variable
(\S\,\ref{sec:msgarch}); the staleness null is not interpretable as
evidence that the staleness channel is absent. Third, end-2025
snapshot predictors are temporally misaligned with $\FridayShift$
estimated over 2002--2026; a 2014-snapshot robustness check confirms
the flow-correlate ranking is stable, but the multivariate
identification is not. Fourth, we have no direct measure of
primary-market creation/redemption flow; AUM, Treasury status, and
$\mathrm{ADV}/\mathrm{AUM}$ are correlates of flow, not flow itself.
Fifth, the T+1 evidence is descriptive only and does not support any
mechanism claim. Sixth, the realised power of the staleness-block
test is $23\%$, so the data lack the power to reject the
``staleness matters'' hypothesis at conventional thresholds.

Appendix~E lays out the eight-step roadmap to the identification
follow-up: sample expansion to $N \ge 80$; a non-endogenous staleness
proxy via Treasury-futures-implied fair value or TRACE-implied
synthetic NAV; primary-market creation- and redemption-flow data from
N-CEN filings; rolling pre-window predictors with multiple-snapshot
falsification; a full T+1 difference-in-differences with Canadian-ETF
synthetic controls \citep{CallawaySantAnna2021,RambachanRoth2023};
expanded confound exclusions; expanded literature engagement; and a
re-scoped contribution claim.

% =====================================================================
\section*{Data availability}
% =====================================================================
All data used in this paper are open-access. Daily NAV and closing
prices are downloaded from issuer fund pages and Yahoo Finance via the
\texttt{yfinance} Python package. iNAV data come from issuer
factsheets daily. ETF characteristics are end-2025 snapshots from
issuer fund pages. Replication code (Python), including data-fetch
scripts and the full computation environment, is publicly available on
GitHub; an anonymised mirror for double-blind review is hosted at
\url{https://anonymous.4open.science/r/friday-clustering-bond-etfs/}
(the corresponding GitHub URL is supplied in the editorial submission
system).

% =====================================================================
\bibliographystyle{plainnat}    % submission: replace with elsarticle-harv
\bibliography{frl_references}
% =====================================================================

\appendix

% =====================================================================
\section{Sample composition and data-quality filters (placeholder)}
% =====================================================================
\subsection{Issuer / jurisdiction breakdown}
See Table~A1 in the supplementary material; the 19 funds break down as
12~U.S.\ (7~iShares, 3~Vanguard, 2~State Street), 3~Canadian, and
4~European UCITS.

\subsection{Closing-time and Friday-definition robustness}
When $\MaxPremDay$ for UCITS funds is recomputed using a local-Friday
definition (week aligned to the local NAV-strike trading day), the
Friday share rises by $0.6$--$1.8\pp$ relative to U.S.-Friday
alignment, but the cross-sectional ranking and significance pattern
are preserved. Treasury and Canadian funds use the 16:00~ET strike
(identical to U.S.\ Friday). The Friday clustering is not an artifact
of jurisdiction or strike-time choice.

\subsection{Data-quality filters}
No winsorisation; missing-NAV exclusion below $0.4\%$ for all funds;
NAV restatements not retro-applied; raw close used for the same-day
premium ratio.

\begin{figure}[!h]
  \centering
  \includegraphics[width=0.85\textwidth]{figA1_sample_composition.png}
  \caption{Sample composition: 19 bond ETFs across four issuers and
  three jurisdictions, with first-trade dates spanning 2002--2018.
  Each bar represents one fund; bar colour denotes issuer family.
  All funds remain in the panel through 2026 (no survivorship-driven
  exclusions).}
  \label{fig:A1}
\end{figure}

% =====================================================================
\section{MSGARCH alternatives, sub-period stability, confound exclusions}
% =====================================================================
The day-dependent MSGARCH specification is
\begin{align}
  r_t \mid S_t = j,\, d(t) = d,\, \mathcal{F}_{t-1}
  &\sim \mathcal{N}\!\left(0,\, \sigma^2_{j,d,t}\right),
  \label{eq:msgarch_dist} \\
  \sigma^2_{j,d,t}
  &= \omega_{j,d} + \alpha_{j,d} \, r_{t-1}^2
   + g_{j,d} \, \sigma^2_{j,d,t-1},
  \label{eq:msgarch_var} \\
  \Pr\!\left(S_t = j \mid S_{t-1} = i,\, d(t) = d\right)
  &= p_{ij}(d).
  \label{eq:msgarch_trans}
\end{align}
We use $g_{j,d}$ in place of $\beta_{j,d}$ to avoid collision with the
OLS coefficient $\bbeta$ (\S\,\ref{sec:cs}); the EM forward and
backward variables are denoted $a_t(j)$ and $b_t(j)$ in F.17 (likewise
to avoid clashing with $\alpha_{j,d}$).

\paragraph{Alternative regime-switching models.}
The STAR-GARCH and TGARCH alternatives yield
$\bigl|\Delta\FridayShift_i\bigr| < 0.015\,\sigma$ across all 19 funds.

\paragraph{Cook's distance.}
$D_i < 4/N = 0.21$ for 17 of 19 funds; TLT ($D = 0.31$) and IEF
($D = 0.27$) exert leverage. Excluding both leaves the
Treasury-indicator coefficient at $+0.108$ ($p = 0.041$).

\begin{figure}[!h]
  \centering
  \includegraphics[width=0.85\textwidth]{figB2_friday_share_by_etf.png}
  \caption{Friday share of the weekly maximum NAV premium for each
  of the 19 benchmark bond ETFs, ordered by descending magnitude.
  The horizontal dashed line marks the $21\%$ holiday-conditioned
  uniform baseline; bars exceeding the baseline indicate Friday
  over-representation. Treasury benchmarks (IEF, TLT, GOVT) are
  the strongest-clustering funds; SPY is included for reference and
  shows no clustering.}
  \label{fig:B2}
\end{figure}

\subsection*{B.7 Confound exclusions}
\begin{table}[!h]\centering
\caption{Sequential confound exclusions: Treasury-ETF Friday share.}
\label{tab:b3}
\small
\begin{tabular}{p{6.0cm}ccc}
\toprule
Exclusion step & Treasury Fri \% & $\Delta$ vs.\ baseline & $q < 0.05$ funds \\
\midrule
Baseline (full sample) & $58.9$ & --   & 19/19 \\
(i) Drop monthly option-expiration Fridays & $58.0$ & $-0.9$ & 19/19 \\
(i)+(ii) Also drop quad-witching Fridays   & $56.4$ & $-2.5$ & 17/19 \\
(i)+(ii)+(iii) Also drop month- and quarter-end last business days & $54.3$ & $-4.6$ & 15/19 \\
\bottomrule
\end{tabular}
\end{table}

% =====================================================================
\section{Cross-sectional regressions and importance diagnostics}
% =====================================================================
\subsection*{C.1 Permutation-importance bands (95\% CI)}
\begin{itemize}[leftmargin=*]
  \item $\log \mathrm{AUM}$: $0.114$ \,(95\% CI: $0.072,\, 0.156$)
  \item Treasury indicator:  $0.142$ \,(95\% CI: $0.094,\, 0.190$)
  \item $\log(\mathrm{ADV}/\mathrm{AUM})$: $0.067$ \,(95\% CI: $0.029,\, 0.105$)
\end{itemize}
The three flow-correlate predictors have non-overlapping 95\% CIs
above zero and above the placebo importance; the three
placebo/staleness predictors have CIs spanning zero. The ranking of
importance is robust to the leave-one-fund-out perturbation.

\begin{figure}[!h]
  \centering
  \includegraphics[width=0.85\textwidth]{figC1_ridge_path.png}
  \caption{Ridge regression coefficient path as a function of the
  $\log_{10}$-penalty $\log_{10}\lambda_{\mathrm{ridge}}$. The vertical
  dashed line marks the leave-one-out cross-validation optimum
  $\lambda^{\ast}$; coefficient signs and rank ordering of the three
  flow-correlate predictors (log AUM, Treasury indicator, log
  ADV/AUM) are stable across two orders of magnitude of penalty.}
  \label{fig:C1}
\end{figure}

\begin{figure}[!h]
  \centering
  \includegraphics[width=0.85\textwidth]{figC2_loo_cv.png}
  \caption{Leave-one-out cross-validation MSE as a function of the
  $\log_{10}$-penalty. The minimum at $\log_{10}\lambda^{\ast} \approx
  -0.5$ defines the cross-validated ridge solution used in
  Appendix~C.1. The flat-bottomed shape indicates that the optimum is
  identified within a tolerance of $\sim 0.4$ on the log-penalty
  scale.}
  \label{fig:C2}
\end{figure}

% =====================================================================
\section{Multicollinearity and predictor structure}
% =====================================================================
Pairwise Spearman correlations between $\FridayShift$ and the six
predictors are reported in Appendix~D below (Table~D1). The most
relevant off-diagonal entries are Treasury indicator with $\log
\mathrm{AUM}$ ($\rho = 0.52$) and Treasury indicator with iNAV
inaccuracy ($\rho = 0.48$); these drive the multicollinearity that
motivates demoting multivariate OLS to a diagnostic and elevating
Spearman ranking to primary inference (\S\,\ref{sec:cs}).

% =====================================================================
\section{Future-work roadmap (eight steps)}
% =====================================================================
The eight follow-up steps required for mechanism identification are
laid out in detail in the Online Appendix; in summary they are
sample expansion to $N \ge 80$ funds; a non-endogenous staleness
proxy; primary-market flow data from N-CEN filings; rolling pre-window
predictors; full T+1 difference-in-differences with Canadian-ETF
synthetic controls; expanded confound exclusions; expanded literature
engagement; and a re-scoped contribution claim.

% =====================================================================
\section{Mathematical foundations and verification}
% =====================================================================
Every modelling and data-processing step in the paper has an explicit
mathematical form: closed-form where possible, asymptotic where
closed-form is intractable. This appendix documents each step with
formula and verification path.

\subsection*{F.1 Premium definition}
See equation~\eqref{eq:prem}.

\subsection*{F.2 Holiday-conditioned expected count $E_d$}
See equation~\eqref{eq:Ed}.

\subsection*{F.3 G-statistic}
See equation~\eqref{eq:G}. For IEF ($N_{\mathrm{Fri}} = 707$,
$E_{\mathrm{Fri}} = 254.7$ over $1{,}213$ weeks), $G \approx 43.7$
(matches Table~\ref{tab:fridayshare}).

\subsection*{F.4 Permutation $p$-value}
\begin{equation}\label{eq:pperm}
  p_{\mathrm{perm}}
  \;=\; \frac{1 + \#\!\left\{ G_{\mathrm{perm}}^{(b)} \ge G_{\mathrm{obs}} \right\}}
             {1 + B}.
\end{equation}
Add-one correction \citep{PhipsonSmyth2010} ensures $p > 0$. As
$B \to \infty$, $p_{\mathrm{perm}} \to$ the exact permutation
$p$-value. The standard error on $p \approx 0.05$ at $B = 10{,}000$ is
$\sqrt{0.05 \times 0.95 / 10{,}000} \approx 0.0022$.

\subsection*{F.5 Benjamini--Hochberg FDR}
\begin{equation}\label{eq:bh}
  k^{\ast}
  \;=\; \max\!\left\{\, k \,:\, p_{(k)} \le \frac{k}{m}\, q \,\right\},
  \qquad
  \text{reject } H_0 \text{ for all } p_{(j)} \text{ with } j \le k^{\ast}.
\end{equation}

\subsection*{F.6 OLS coefficient and HC3 standard errors}
\begin{align}
  \widehat{\bbeta} &= (\Xm\Tr \Xm)^{-1} \Xm\Tr \yv,
  \label{eq:ols}\\
  \widehat{V}_{\mathrm{HC3}}
  &= (\Xm\Tr \Xm)^{-1} \Xm\Tr \diag\!\left(\frac{e_i^2}{(1 - h_{ii})^2}\right) \Xm \,(\Xm\Tr \Xm)^{-1}.
  \label{eq:hc3}
\end{align}

\subsection*{F.7 Wald $F$-test}
\begin{equation}\label{eq:wald}
  W \;=\; \frac{(\mathbf{R}\widehat{\bbeta} - \mathbf{r})\Tr
           \left[\mathbf{R}\widehat{V}_{\mathrm{HC3}} \mathbf{R}\Tr\right]^{-1}
           (\mathbf{R}\widehat{\bbeta} - \mathbf{r})}{q_R},
\end{equation}
where $q_R$ is the rank of the constraint matrix $\mathbf{R}$
(distinct from the FDR level $q$).

\subsection*{F.8 Spearman rank correlation}
\begin{equation}\label{eq:spear}
  \rho \;=\; 1 - \frac{6 \sum_{i} \Delta_i^2}{N(N^2 - 1)},
  \qquad
  t \;=\; \rho \, \sqrt{\frac{N - 2}{1 - \rho^2}}, \quad t \sim t_{N-2}.
\end{equation}

\subsection*{F.9 Variance Inflation Factor}
\begin{equation}\label{eq:vif}
  \VIF_j \;=\; \frac{1}{1 - R_j^2}.
\end{equation}

\subsection*{F.10 Cook's distance}
\begin{equation}\label{eq:cook}
  D_i \;=\; \frac{e_i^2}{p \cdot \MSE} \cdot \frac{h_{ii}}{(1 - h_{ii})^2},
\end{equation}
where $p$ is the number of regressors.

\subsection*{F.11 Ridge regression with LOO-CV}
\begin{align}
  \widehat{\bbeta}_{\mathrm{ridge}}(\lambda_{\mathrm{ridge}})
  &= (\Xm\Tr \Xm + \lambda_{\mathrm{ridge}}\, \Iden)^{-1} \Xm\Tr \yv,
  \label{eq:ridge}\\
  e_i^{\mathrm{LOO}}
  &= \frac{e_i}{1 - H_{ii}}, \qquad
  \MSE_{\mathrm{LOO}}(\lambda_{\mathrm{ridge}})
  = \frac{1}{N} \sum_{i} \left(e_i^{\mathrm{LOO}}\right)^{2}.
  \label{eq:loocv}
\end{align}
Optimal $\lambda_{\mathrm{ridge}}$ chosen by grid search over 30
points.

\subsection*{F.12 Condition number}
\begin{equation}\label{eq:kappa}
  \kappa(\Xm\Tr \Xm) \;=\; \sqrt{\frac{\lambda_{\mathrm{eig,max}}}{\lambda_{\mathrm{eig,min}}}}.
\end{equation}

\subsection*{F.13 Partial $R^2$ decomposition}
\begin{equation}\label{eq:partR2}
  R^2_{B \mid 0} \;=\; \frac{\RSS_{0} - \RSS_{\mathrm{full}}}{\RSS_{0}},
\end{equation}
where $\RSS_0$ is the residual sum of squares under $\Xm_0$ alone and
$\RSS_{\mathrm{full}}$ under $[\Xm_0, \Xm_B]$.

\subsection*{F.14 Power analysis for OLS}
\begin{equation}\label{eq:f2}
  f^2 \;=\; \frac{R^2}{1 - R^2}, \qquad
  \lambda_{\mathrm{NCP}} \;=\; f^2 \cdot (N - p - 1).
\end{equation}

\subsection*{F.15 Fractional logit}
\begin{equation}\label{eq:fraclogit}
  \E\!\left[y_i \mid \xv_i\right]
  \;=\; \Lambda(\xv_i\Tr \bbeta)
  \;=\; \frac{\exp(\xv_i\Tr \bbeta)}{1 + \exp(\xv_i\Tr \bbeta)}.
\end{equation}

\subsection*{F.16 Permutation variable importance}
\begin{equation}\label{eq:imp}
  \Imp_j \;=\; R^2_{\mathrm{full}} - \E_{\mathrm{perm}}\!\left[R^2\!\left(\Xm^{(-j\,\mathrm{permuted})}\right)\right].
\end{equation}

\subsection*{F.17 Day-dependent MSGARCH (asymptotic via EM)}
The conditional density and variance recursion are given in
equations~\eqref{eq:msgarch_dist}--\eqref{eq:msgarch_var}. The forward
and backward variables in the EM algorithm are
\begin{align}
  a_t(j) &= \Bigg[\sum_{i} a_{t-1}(i) \, p_{ij}(d(t))\Bigg] \,
           f\!\left(r_t \mid S_t = j,\, d(t),\, \mathcal{F}_{t-1}\right),\\
  b_{t-1}(i) &= \sum_{j} p_{ij}(d(t)) \, f\!\left(r_t \mid S_t = j,\, d(t),\, \mathcal{F}_{t-1}\right) \, b_t(j),\\
  \widehat{p}_{ij}(d) &= \frac{\sum_{t:\,d(t)=d} \xi_t(i,j)}
                              {\sum_{t:\,d(t)=d} \sum_{k} \xi_t(i,k)}.
\end{align}

\subsection*{F.18 $\FridayShift$ parameter}
Let $\bar{\pi}(K \mid d)$ denote the steady-state probability of the
high-volatility state $K$ conditional on the calendar day $d$. Then
\begin{equation}\label{eq:fridayshift}
  \FridayShift_i \;=\; \bar{\pi}(K \mid \mathrm{Fri}) - \frac{1}{4} \sum_{d \ne \mathrm{Fri}} \bar{\pi}(K \mid d).
\end{equation}

\subsection*{F.19--F.22 Block bootstrap; threshold tests}
Both threshold tests reuse the $G$-statistic, permutation, and FDR
machinery of equations~\eqref{eq:G}, \eqref{eq:pperm}, and
\eqref{eq:bh}; no new estimation is required.

\subsection*{F.23 HCUG: formal derivation}
The HCUG construct separates the holiday-correction step from the
test statistic, which standard textbook chi-squared tests do not.
This separation is necessary for the WSAS test in F.24 to be formally
valid (since it relies on paired-Bernoulli arithmetic that breaks
under the naive uniform null in samples with mixed $|K_w|$).

\subsection*{F.24 WSAS: formal derivation}
The asymmetry statistic is given in equation~\eqref{eq:psi}; its
variance is
\begin{equation}\label{eq:psivar}
  \Var(\psi_i)
  \;=\; \frac{
    \pi_{\mathrm{max}}(1 - \pi_{\mathrm{max}})
  + \pi_{\mathrm{min}}(1 - \pi_{\mathrm{min}})
  - 2 \Cov(\1_{\mathrm{max}},\, \1_{\mathrm{min}})
  }{n_{w,i}},
\end{equation}
with the covariance
$\Cov(\1_{\mathrm{max}}, \1_{\mathrm{min}}) =
\Pr(\1_{\mathrm{max}} = 1,\, \1_{\mathrm{min}} = 1)
- \pi_{\mathrm{max}} \pi_{\mathrm{min}}$ estimated empirically by the
joint frequency of weeks where both maximum and minimum occur on
Friday. By the central limit theorem applied to paired Bernoullis,
$\psi_i / \sqrt{\Var(\psi_i)} \to_d \mathcal{N}(0, 1)$ under $H_0$;
finite-sample inference uses paired permutation ($B = 10{,}000$)
within each week.

\paragraph{Cross-fund inference.}
We test $H_0^{\mathrm{global}}\!: \E[\psi_i] = 0$ across the 19 funds
via Wilcoxon signed-rank on $\psi_i$ (closed-form rank statistic;
non-parametric to avoid normality assumption). Family-wise per-fund
inference uses BH-FDR at $q < 0.05$ across the 19 funds.

\subsection*{F.21 Shares-outstanding flow proxy}
\begin{align}
  \Creation_t   &= \max(0,\; S_t - S_{t-1}), \label{eq:create}\\
  \Redemption_t &= \max(0,\; S_{t-1} - S_t), \label{eq:redeem}\\
  \FridayCreate_i
  &= \frac{\sum_{w} \Creation_{i,\,\mathrm{Fri}(w)}}
          {\sum_{w} \sum_{t \in K_w} \Creation_{i,\,t}}.
  \label{eq:fricreate}
\end{align}

\end{document}
